% © 2026 Ing. David Jaroš — CC BY-NC-ND 4.0
%
% This work is licensed under a Creative Commons Attribution-NonCommercial-NoDerivatives
% 4.0 International License (CC BY-NC-ND 4.0).
%
% License History: Earlier drafts (up to v0.3) were released under CC BY 4.0.
% From v0.4 onward, all material is released under CC BY-NC-ND 4.0 to protect
% the integrity of the theoretical work during ongoing academic development.
%
% See LICENSE.md for full license text.
%
% File: research_tracks/alpha_spectral/hecke_equivariant_path_integral.tex
%
% Task: construct_or_kill_hecke_equivariant_winding_path_integral
% Priority: CRITICAL
% Date: 2026-05-09
%
% Purpose:
%   Determine whether the UBT winding sector at prime level p admits a
%   Hecke-equivariant path-integral decomposition whose saddle count equals
%   mu(Gamma_0(p)) = p+1.  Constructs each required ingredient in turn and
%   identifies the precise obstruction when construction fails.
%
% Hard rules:
%   - No observed alpha as input anywhere.
%   - No fitting to 137.
%   - No speculative consciousness or dark-sector interpretation.
%   - If construction fails, state the exact obstruction.

% UBT-AI-PROVENANCE-BEGIN
% schema: ubt-ai-provenance/v1
% tier: C_working
% ai_assistance: disclosed
% human_review: risk-based
% editorial_responsibility: Ing. David Jaroš
% policy: ../../AI_PROVENANCE.md
% notice: Working material; exhaustive human review is not claimed.
% UBT-AI-PROVENANCE-END

\documentclass[12pt,a4paper]{article}
\usepackage[a4paper, margin=2.5cm]{geometry}
\usepackage{amsmath,amssymb,amsthm}
\usepackage{mathtools}
\usepackage{hyperref}
\usepackage{xcolor}
\usepackage{booktabs}
\usepackage{longtable}
\usepackage{mdframed}
\usepackage{tcolorbox}
\tcbuselibrary{skins,breakable}
\usepackage{array}
\usepackage{enumitem}

%──────────────────────────────── theorem environments ─────────────────────────
\newtheorem{theorem}{Theorem}[section]
\newtheorem{lemma}[theorem]{Lemma}
\newtheorem{proposition}[theorem]{Proposition}
\newtheorem{corollary}[theorem]{Corollary}
\theoremstyle{definition}
\newtheorem{definition}[theorem]{Definition}
\newtheorem{construction}[theorem]{Construction Attempt}
\theoremstyle{remark}
\newtheorem{remark}[theorem]{Remark}

%──────────────────────────────── coloured boxes ────────────────────────────────
\newmdenv[backgroundcolor=green!8, linecolor=green!50!black, linewidth=1.2pt]{resultbox}
\newmdenv[backgroundcolor=yellow!12,linecolor=orange!80!black,linewidth=1.2pt]{gapbox}
\newmdenv[backgroundcolor=red!7,   linecolor=red!55!black,   linewidth=1.2pt]{warnbox}
\newmdenv[backgroundcolor=blue!6,  linecolor=blue!55!black,  linewidth=1.2pt]{insightbox}

\newcommand{\PROVED}{\textcolor{green!50!black}{\textbf{[PROVED]}}}
\newcommand{\OPEN}  {\textcolor{red!70!black}  {\textbf{[OPEN]}}}
\newcommand{\COND}  {\textcolor{blue!70!black} {\textbf{[CONDITIONAL]}}}
\newcommand{\NOGO}  {\textcolor{red!80!black}  {\textbf{[NO-GO]}}}
\newcommand{\OBS}   {\textcolor{orange!85!black}{\textbf{[OBSTRUCTION]}}}

%──────────────────────────────── macros ────────────────────────────────────────
\newcommand{\SL}{\mathrm{SL}}
\newcommand{\PSL}{\mathrm{PSL}}
\newcommand{\Tr}{\mathrm{Tr}}
\newcommand{\vol}{\operatorname{vol}}
\newcommand{\HH}{\mathbb{H}}       % quaternion algebra
\newcommand{\bH}{\mathfrak{H}}     % upper half-plane
\newcommand{\CC}{\mathbb{C}}
\newcommand{\RR}{\mathbb{R}}
\newcommand{\ZZ}{\mathbb{Z}}
\newcommand{\FF}{\mathbb{F}}
\newcommand{\PP}{\mathbb{P}}
\newcommand{\Gam}{\Gamma}
\newcommand{\bTheta}{\bar{\Theta}}
\newcommand{\Hp}{\mathcal{H}_p}
\newcommand{\Zp}{Z_p}

%──────────────────────────────────────────────────────────────────────────────
\title{\textbf{Hecke-Equivariant Path-Integral Decomposition\\
       of the UBT Winding Sector at Prime Level $p$}\\[4pt]
       \large Construction Attempt and No-Go Analysis}
\author{Ing.\ David Jaro\v{s}}
\date{2026-05-09}

\begin{document}
\maketitle

\begin{abstract}
We carry out a systematic construction of the Hecke-equivariant
path-integral decomposition of the Unified Biquaternion Theory (UBT)
winding sector at prime level $p$.  The required ingredients are:
a Hilbert space of winding states $\Hp$; an action of $\SL(2,\ZZ)$ or
Hecke correspondences on $\Hp$; a quotient identification
$\Gam_0(p)\backslash\SL(2,\ZZ)\cong\PP^1(\FF_p)$; a candidate saddle set
indexed by $\PP^1(\FF_p)$; one-loop weights for each saddle; and a
derivation of the normalisation by
$\vol(\SL(2,\ZZ)\backslash\bH)=\pi/3$.
Each step is carried to the point at which it either succeeds or hits an
obstruction.  Result: the arithmetic layer (coset count $p+1$, modular
area, normalisation by $\pi/3$) is exact.  The physics layer (modular
$\tau$ variable from $S[\Theta]$, $\Gam_0(p)$ equivariance of the
measure, equal-action saddles) encounters three sharp obstructions that
prevent a first-principles derivation.  The construction is
\textbf{incomplete}: no go at the level of the current UBT action, with
exact obstruction statements.
\end{abstract}

\begin{warnbox}
\textbf{Hard constraint.} No observed value of $\alpha$ is used at any
point.  No fitting to 137.  No consciousness or dark-sector
interpretation.  If construction fails, the exact obstruction is
recorded.
\end{warnbox}

\tableofcontents
\newpage

%──────────────────────────────────────────────────────────────────────────────
\section{Problem statement and notation}
\label{sec:problem}
%──────────────────────────────────────────────────────────────────────────────

The effective potential used in the T3\_ALPHA route is
\begin{equation}
  V_{\mathrm{eff}}(n) = n^2 - B\,n\ln n, \qquad n\in\ZZ_{>0}.
  \label{eq:veff}
\end{equation}
The stationary condition $\partial V_{\mathrm{eff}}/\partial n=0$ at
$n^*=p$ (a prime) gives
\begin{equation}
  2p = B(\ln p + 1).
  \label{eq:stat}
\end{equation}
The modular ansatz proposes
\begin{equation}
  B(p) = \frac{\mu(\Gam_0(p))}{3} = \frac{p+1}{3},
  \label{eq:Bansatz}
\end{equation}
where $\mu(\Gam_0(p))=[\SL(2,\ZZ):\Gam_0(p)]$ is the index of the
congruence subgroup
\[
  \Gam_0(p) = \left\{\begin{pmatrix}a&b\\c&d\end{pmatrix}\in\SL(2,\ZZ)
              \;\middle|\; c\equiv 0\pmod{p}\right\}.
\]
The goal is to determine whether~\eqref{eq:Bansatz} can be derived from
the UBT action $S[\Theta]$ via a Hecke-equivariant path-integral
decomposition, or whether a precise obstruction prevents this.

\paragraph{Questions to be answered.}
\begin{enumerate}[label=Q\arabic*.]
  \item Does $S[\Theta]$ naturally define a modular $\tau$ variable?
  \item Does the $p$-winding sector transform under $\Gam_0(p)$?
  \item Are there exactly $p+1$ inequivalent saddles?
  \item Are they equal-action?
  \item Does their contribution multiply the $n\ln n$ coefficient?
  \item Does the division by $3$ enter through modular area normalisation?
\end{enumerate}

%──────────────────────────────────────────────────────────────────────────────
\section{Arithmetic layer: exact results}
\label{sec:arithmetic}
%──────────────────────────────────────────────────────────────────────────────

Before any physics enters, the following identities are exact and
unconditional.

\subsection{Coset count}

\begin{lemma}[Index of $\Gam_0(p)$]
\label{lem:index}
For prime $p$,
\[
  \mu(\Gam_0(p)) = [\SL(2,\ZZ):\Gam_0(p)] = p+1.
\]
\end{lemma}

\begin{proof}
By the index formula for congruence subgroups,
$\mu(\Gam_0(N))=N\prod_{q|N}(1+q^{-1})$.  For $N=p$ prime this gives
$p(1+1/p)=p+1$.
\end{proof}

\begin{lemma}[Coset representatives via $\PP^1(\FF_p)$]
\label{lem:cosets}
There is a bijection
\[
  \Gam_0(p)\backslash\SL(2,\ZZ)
  \;\longleftrightarrow\;
  \PP^1(\FF_p) = \{[0:1],[1:0],[1:1],\ldots,[p-1:1]\},
\]
so the right-coset space has cardinality $p+1$.
\end{lemma}

\begin{proof}
Map $\left(\begin{smallmatrix}a&b\\c&d\end{smallmatrix}\right)$ to
$[c:d]\in\PP^1(\FF_p)$.  This is well-defined and $\Gam_0(p)$-equivariant
on the left.  The map is surjective with trivial kernel in the quotient.
\end{proof}

\subsection{Modular area normalisation}

\begin{lemma}[Fundamental-domain volume]
\label{lem:fundvol}
\[
  \vol\!\bigl(\SL(2,\ZZ)\backslash\bH\bigr) = \frac{\pi}{3}.
\]
\end{lemma}

\begin{proof}
The standard fundamental domain $\mathcal{F}$ for $\SL(2,\ZZ)$ on the
upper half-plane $\bH$ has hyperbolic area $\pi/3$ (see, e.g.,
\cite{Serre:1973}).
\end{proof}

\begin{corollary}[Area of $X_0(p)$]
\label{cor:area}
\[
  \vol(X_0(p)) = \frac{\pi}{3}\,\mu(\Gam_0(p)) = \frac{\pi(p+1)}{3},
  \qquad
  \frac{\vol(X_0(p))}{\pi} = \frac{p+1}{3} = B(p).
\]
\end{corollary}

\begin{resultbox}
\textbf{Arithmetic summary.}  The count $p+1$ is exact (Lemma~\ref{lem:index}).
The bijection with $\PP^1(\FF_p)$ is exact (Lemma~\ref{lem:cosets}).
The normalisation to $(p+1)/3$ via modular area is exact
(Corollary~\ref{cor:area}).  These results hold independently of any
physics.
\end{resultbox}

%──────────────────────────────────────────────────────────────────────────────
\section{Physics layer: constructing $\mathcal{H}_p$}
\label{sec:hilbert}
%──────────────────────────────────────────────────────────────────────────────

\subsection{UBT action and complex time}

The UBT action on a background $g_{\mu\nu}$ and complex time
$\tau = t + i\psi$ is taken to be
\begin{equation}
  S[\Theta] = \int_{\mathcal{M}}
  \bigl(\nabla^\dagger\Theta\bigr)^\dagger
  \bigl(\nabla^\dagger\Theta\bigr)\,
  \sqrt{-g}\,d^4x\,d\psi,
  \label{eq:action}
\end{equation}
where $\Theta(q,\tau)\in\CC\otimes\HH$ is the fundamental biquaternion
field and $\nabla^\dagger$ is the biquaternionic Dirac--Laplace operator.

\subsection{Winding modes in the $\psi$-direction}

The imaginary-time direction $\psi$ is treated as compact with period
$\beta$; the Fourier expansion is
\begin{equation}
  \Theta(x,t,\psi) = \sum_{n\in\ZZ} \Theta_n(x,t)\,e^{2\pi i n\psi/\beta}.
  \label{eq:winding}
\end{equation}
The mode $\Theta_n$ carries \emph{winding number} $n$ (Kaluza--Klein
charge in the $\psi$-direction).

\begin{definition}[Winding-sector Hilbert space $\Hp$]
\label{def:Hp}
The \emph{level-$p$ winding sector} at prime $p$ is
\[
  \Hp = \overline{\mathrm{span}}\,\{\Theta_n \mid n\equiv 0\pmod{p}\},
\]
the $L^2$ closure of the subspace of winding modes whose winding number
is divisible by $p$.
\end{definition}

\begin{remark}
$\Hp$ is a natural subspace because modes with $n\equiv 0\pmod{p}$ close
under the quadratic part of $S[\Theta]$: the kinetic term does not mix
$n\equiv 0$ and $n\not\equiv 0$ modes on a symmetric background.
\end{remark}

%──────────────────────────────────────────────────────────────────────────────
\section{Construction attempt 1: $\SL(2,\ZZ)$ action on $\Hp$}
\label{sec:sl2z}
%──────────────────────────────────────────────────────────────────────────────

\subsection{Candidate action}

The upper half-plane variable is identified with the complex-time torus
modulus
\begin{equation}
  \tau = \frac{t + i\psi}{\beta} \in \bH.
  \label{eq:tau}
\end{equation}
The group $\SL(2,\ZZ)$ acts on $\tau$ by Möbius transformations
$\gamma\colon\tau\mapsto(a\tau+b)/(c\tau+d)$ for
$\gamma=\left(\begin{smallmatrix}a&b\\c&d\end{smallmatrix}\right)$.

A candidate representation of $\SL(2,\ZZ)$ on $\Hp$ is defined by
requiring that the path-integral measure $\mathcal{D}\Theta_n$ transforms
as a weight-$(-\tfrac{1}{2})$ modular form in each mode, so that the
partition function
\[
  Z(\tau) = \int_{\Hp} \mathcal{D}\Theta_p\,e^{-S[\Theta]}
\]
is modular invariant (weight 0) under $\SL(2,\ZZ)$.

\subsection{Obstruction 1: modular invariance of $S[\Theta]$}

\begin{proposition}[Obstruction O1]
\label{prop:O1}
The UBT action~\eqref{eq:action} is \textbf{not} demonstrated to be
modular invariant under $\SL(2,\ZZ)$ acting by $\tau\mapsto(a\tau+b)/(c\tau+d)$.
\end{proposition}

\begin{proof}
The action $S[\Theta]$ is an integral over a fixed spacetime manifold
$\mathcal{M}$ with the complex-time $\psi$ direction treated as a compact
extra dimension.  Under $\tau\to(a\tau+b)/(c\tau+d)$ the ratio
$t/\beta$ and $\psi/\beta$ mix.  For this to be a symmetry of
$S[\Theta]$, the kinetic operator $\nabla^\dagger\nabla$ would need to
transform covariantly under reparametrisations of the complex-time torus,
i.e., $\nabla^\dagger\nabla$ would need to be a modular covariant operator
of definite weight.

The current derivations (see \texttt{canonical/fields/}) establish the
form of $\nabla^\dagger\nabla$ in a fixed background metric but do not
show that it transforms as a modular form under $\SL(2,\ZZ)$.  In
particular, the integration measure $d^4x\,d\psi$ picks up a Jacobian
$|c\tau+d|^{-2}$ under $\tau\to(a\tau+b)/(c\tau+d)$, which is not
cancelled by the transformation of the kinetic term unless additional
modular-weight cancellations are proved.  No such proof exists in the
repository.
\end{proof}

\begin{gapbox}
\textbf{Obstruction O1.} The action $S[\Theta]$ does not define a modular
$\tau$ variable in the sense of being $\SL(2,\ZZ)$-invariant.  Proving
$\SL(2,\ZZ)$ symmetry would require showing that
$\nabla^\dagger\nabla$ on the complex-time torus is a modular covariant
operator, which is currently unestablished.
\end{gapbox}

%──────────────────────────────────────────────────────────────────────────────
\section{Construction attempt 2: $\Gam_0(p)$ equivariance of $\Hp$}
\label{sec:gamma0}
%──────────────────────────────────────────────────────────────────────────────

\subsection{Arithmetic compatibility}

Suppose, conditionally on O1 being resolved, that $S[\Theta]$ is
$\SL(2,\ZZ)$-invariant.  Then one can ask whether the subspace $\Hp$
(modes with $n\equiv 0\pmod p$) is preserved by $\Gam_0(p)$.

\begin{lemma}[Arithmetic compatibility]
\label{lem:arith}
Conditional on $\SL(2,\ZZ)$ acting on winding modes, the subspace $\Hp$
of modes with $n\equiv 0\pmod{p}$ is preserved by $\Gam_0(p)$.
\end{lemma}

\begin{proof}
Under a torus reparametrisation by
$\gamma=\left(\begin{smallmatrix}a&b\\c&d\end{smallmatrix}\right)
\in\Gam_0(p)$, the new winding number of a mode originally at $n$ is
$n'=cn+d\cdot(\text{topological term})$; schematically the sublattice
$p\ZZ\subset\ZZ$ is preserved under the action of $\Gam_0(p)$ because
$c\equiv 0\pmod{p}$ implies $cn\equiv 0\pmod{p}$.  Hence $\Hp$ is
closed under $\Gam_0(p)$.
\end{proof}

\subsection{Obstruction 2: the $\SL(2,\ZZ)$-action on winding numbers
is not derived}

\begin{proposition}[Obstruction O2]
\label{prop:O2}
No derivation from $S[\Theta]$ establishes that $\SL(2,\ZZ)$ acts on
winding modes $\Theta_n$ by a transformation that maps the winding-number
label $n$ according to a group homomorphism $\SL(2,\ZZ)\to\mathrm{Aut}(\ZZ)$.
\end{proposition}

\begin{proof}
The winding number $n$ enters $S[\Theta]$ through the kinetic term
$\sim n^2/\beta^2$ for each mode $\Theta_n$.  A Möbius transformation
$\tau\to(a\tau+b)/(c\tau+d)$ rescales the modulus $\beta\to|c\tau+d|\beta$
and mixes real and imaginary components of $\tau$.  For the winding label
$n$ to transform simply (e.g., $n\to cn+d$, or $n\to n/(c\tau+d)$ in some
sense), the Fourier expansion~\eqref{eq:winding} must be reformulated in
terms of the new period.  But the canonical Fourier expansion depends on
the period $\beta$, and rewriting in terms of the transformed period
reshuffles all modes.  There is no derivation showing that this reshuffling
has a well-defined, finite-dimensional action on the subspace $\Hp$.
\end{proof}

\begin{gapbox}
\textbf{Obstruction O2.} Even if $S[\Theta]$ is $\SL(2,\ZZ)$-invariant,
the induced action on winding modes has not been computed.  Lemma~\ref{lem:arith}
holds purely arithmetically, but connecting it to a physical action of
$\Gam_0(p)$ on $\Hp$ requires the missing derivation.
\end{gapbox}

%──────────────────────────────────────────────────────────────────────────────
\section{Construction attempt 3: saddle set indexed by $\PP^1(\FF_p)$}
\label{sec:saddles}
%──────────────────────────────────────────────────────────────────────────────

\subsection{Candidate saddle configurations}

Suppose conditionally that $\Gam_0(p)\backslash\SL(2,\ZZ)$ labels
inequivalent boundary conditions for the path integral over $\Hp$.
For each coset representative
$\gamma_a = \left(\begin{smallmatrix}1&a\\0&1\end{smallmatrix}\right)$
($a=0,1,\ldots,p-1$) and
$\gamma_\infty = \left(\begin{smallmatrix}0&-1\\1&0\end{smallmatrix}\right)$,
define the \emph{$a$-th saddle configuration} as the classical solution
$\bar\Theta_a$ of the equations of motion of $S[\Theta]$ with the torus
modulus shifted to $\gamma_a\cdot\tau$.

\begin{definition}[Candidate saddle set]
\label{def:saddles}
The candidate set of $p+1$ saddles is
\[
  \mathcal{S} = \bigl\{\bar\Theta_a \;\big|\; [a]\in\PP^1(\FF_p)\bigr\},
\]
where each $\bar\Theta_a$ is the classical field at the modular image
$\gamma_a\cdot\tau$ of the reference saddle $\bar\Theta_0$.
\end{definition}

\subsection{One-loop weights}

The one-loop partition function around each saddle is
\begin{equation}
  Z_a^{(1)} = e^{-S[\bar\Theta_a]}
              \Bigl(\det'\nabla^\dagger\nabla\Big|_{\bar\Theta_a}\Bigr)^{-1/2}.
  \label{eq:1loop}
\end{equation}
If $\bar\Theta_a = \gamma_a^*\bar\Theta_0$ (pull-back under the modular
transformation), then the classical action satisfies
$S[\bar\Theta_a]=S[\bar\Theta_0]$ if and only if $S[\Theta]$ is
$\SL(2,\ZZ)$-invariant.  Similarly, the functional determinant satisfies
$\det'\nabla^\dagger\nabla|_{\bar\Theta_a}=\det'\nabla^\dagger\nabla|_{\bar\Theta_0}$
under the same assumption.

\subsection{Condition for equal weighting}

\begin{proposition}[Equal-action condition]
\label{prop:equal}
All $p+1$ candidate saddles carry equal classical action and equal one-loop
determinant if and only if $S[\Theta]$ is invariant under the full group
$\SL(2,\ZZ)$ acting by $\tau$-reparametrisations.
\end{proposition}

\begin{proof}
By definition, the saddles $\bar\Theta_a$ are related by modular
transformations $\gamma_a\in\SL(2,\ZZ)$.  Equal action requires
$S[\gamma_a^*\Theta]=S[\Theta]$ for all $\gamma_a$, i.e., $\SL(2,\ZZ)$
invariance of $S$.  Equal one-loop weights require the spectrum of
$\nabla^\dagger\nabla$ to be invariant, which likewise requires $\SL(2,\ZZ)$
equivariance of the operator.
\end{proof}

\subsection{Obstruction 3: equal-action cannot be established}

\begin{proposition}[Obstruction O3]
\label{prop:O3}
The claim that all $p+1$ candidate saddles~\eqref{def:saddles} are
equal-action cannot be derived from the current UBT framework.
\end{proposition}

\begin{proof}
By Proposition~\ref{prop:equal}, equal action requires $\SL(2,\ZZ)$
invariance of $S[\Theta]$.  By Obstruction O1
(Proposition~\ref{prop:O1}), this invariance has not been proved.
Therefore equal weighting cannot be established.  Conversely, without
equal weighting the saddle sum
\[
  \Zp = \sum_{a\in\PP^1(\FF_p)} Z_a^{(1)}
\]
does not factor as $(p+1)\cdot e^{-S_*}\cdot(\text{1-loop})$, and the
degeneracy factor $p+1$ does not emerge.
\end{proof}

\begin{gapbox}
\textbf{Obstruction O3.} The equal-action property of the $p+1$ candidate
saddles cannot be deduced without first proving $\SL(2,\ZZ)$ invariance
of $S[\Theta]$, which is Obstruction O1.
\end{gapbox}

%──────────────────────────────────────────────────────────────────────────────
\section{Construction attempt 4: from saddle sum to the $n\ln n$ coefficient}
\label{sec:nlogn}
%──────────────────────────────────────────────────────────────────────────────

\subsection{Required connection}

Suppose, conditionally, that equal-action holds and that
$\Zp = (p+1)e^{-S_*}\cdot C_{\mathrm{1-loop}}$.
To obtain $V_{\mathrm{eff}}(n)=n^2-B\,n\ln n$, one would need:

\begin{enumerate}
  \item The free energy $F_p = -\ln\Zp$ to contribute a term
        $-(p+1)\ln(e^{-S_*}\cdot C_{\mathrm{1-loop}})$ to the effective action
        at winding number $n=p$.
  \item This contribution to reproduce the structure $-B\,n\ln n$ in~\eqref{eq:veff},
        meaning $C_{\mathrm{1-loop}}$ must produce a $\ln p$ term, consistent
        with a one-loop KK determinant $\sim n^{-c}$ on a torus of radius $\beta$.
  \item The coefficient to be exactly $B=(p+1)/3$ after dividing by the modular
        normalisation.
\end{enumerate}

\begin{proposition}[Conditional $n\ln n$ coefficient]
\label{prop:nlogn}
Conditional on O1, O2, and O3 being resolved:
if $C_{\mathrm{1-loop}}=(\pi^2/\beta^2)^{\lambda}$ for some exponent $\lambda$,
then $\ln\Zp$ contributes a term of the form $\const\cdot(p+1)\ln(p/\beta)$,
which matches the structure $B\cdot n\ln n$ for $n=p$ with
$B\propto p+1$ up to a numerical factor depending on $\lambda$.
The prefactor becomes $B=(p+1)/3$ only after the area normalisation of
Section~\ref{sec:modular_norm}.
\end{proposition}

\begin{proof}
The one-loop determinant of $\nabla^\dagger\nabla$ on the level-$p$
winding sector at modulus $\tau$ is schematically
$(\det'\nabla^\dagger\nabla)^{-1/2}\sim e^{c\,\mathrm{Im}(\tau)\cdot p^2/\beta^2}
 \prod_{n\neq 0}(1-e^{2\pi i n\tau})^{-1}$.
In the large-$p$ limit at fixed $\mathrm{Im}(\tau)$, the leading-log
contribution to $-\ln\det'$ is $\sim p^2\cdot(\text{area})$, plus a
$\ln p$ term from the zero-mode counting.  Multiplying by the saddle
degeneracy $(p+1)$ gives the $\sim(p+1)\ln p$ structure.  Setting
$n=p$ reproduces $B\cdot p\ln p$ if $B=p+1$ up to normalisation.  The
normalisation to $(p+1)/3$ is carried out in the next section.
\end{proof}

\begin{insightbox}
The $n\ln n$ structure is consistent with the expected one-loop determinant
on a torus: this aspect is arithmetically coherent.  The barrier is not
the $\ln n$ term itself but the unproved $\SL(2,\ZZ)$ invariance that
would justify the degeneracy factor $p+1$.
\end{insightbox}

%──────────────────────────────────────────────────────────────────────────────
\section{Normalisation by $\vol(\SL(2,\ZZ)\backslash\bH)=\pi/3$}
\label{sec:modular_norm}
%──────────────────────────────────────────────────────────────────────────────

\subsection{Why division by $3$ enters through modular area}

If the path integral is normalised over the modular fundamental domain
$\mathcal{F}=\SL(2,\ZZ)\backslash\bH$ with the hyperbolic measure
$d\mu=d(\mathrm{Re}\,\tau)\,d(\mathrm{Im}\,\tau)/(\mathrm{Im}\,\tau)^2$,
then the total volume is $\vol(\mathcal{F})=\pi/3$ (Lemma~\ref{lem:fundvol}).
The physically natural normalisation of the partition function is then
\begin{equation}
  \hat{Z}(\tau) = \frac{Z(\tau)}{\vol(\mathcal{F})} = \frac{Z(\tau)}{\pi/3} = \frac{3}{\pi}\,Z(\tau).
  \label{eq:normZ}
\end{equation}
Applied to the saddle sum, if $\Zp=(p+1)e^{-S_*}\cdot C$, the
normalised partition function per unit modular area is
\begin{equation}
  \frac{\Zp}{\pi} = \frac{(p+1)\,e^{-S_*}\,C}{\pi}
  = \frac{p+1}{3}\cdot\frac{3\,e^{-S_*}\,C}{\pi}.
  \label{eq:normsaddle}
\end{equation}
The factor $(p+1)/3$ appears as a ratio of modular areas:
\begin{equation}
  \frac{p+1}{3}
  = \frac{\mu(\Gam_0(p))\,\vol(\mathcal{F})}{\pi}
  = \frac{\vol(X_0(p))}{\pi}.
  \label{eq:ratio}
\end{equation}

\begin{resultbox}
\textbf{Exact result.}  The division by $3$ in $B(p)=(p+1)/3$ is
exactly explained as the ratio of the level-$p$ modular curve area to
$\pi$, i.e., the modular-area normalisation of
$\vol(\SL(2,\ZZ)\backslash\bH)=\pi/3$.  This is an unconditional
arithmetic fact (Corollary~\ref{cor:area}).  It does \emph{not} require
any UBT-specific assumption.
\end{resultbox}

\subsection{Alternative: $\mathrm{Im}(\HH)$ / theta-dimension origin}

Internal notes have suggested that the denominator $3$ could alternatively
arise from the three-dimensional imaginary quaternion subspace
$\mathrm{Im}(\HH)\cong\RR^3$.  This would give the mnemonic
\[
  B(p) = \frac{p+1}{\dim_\RR\mathrm{Im}(\HH)}.
\]
However, no one-loop determinant or mode-counting computation in the
current repository derives this identification from $S[\Theta]$.  It is
noted as a structural analogy only.

\begin{proposition}
\label{prop:im3}
The identification of the denominator $3$ with
$\dim_\RR\mathrm{Im}(\HH)=3$ is a heuristic mnemonic.  It is not
stronger than the modular-area derivation.
\end{proposition}

\begin{proof}
The modular-area derivation (Corollary~\ref{cor:area}) is exact and
requires only the definition of $\Gam_0(p)$ and the volume of the
$\SL(2,\ZZ)$ fundamental domain.  The $\mathrm{Im}(\HH)$ identification
would require showing that the one-loop determinant of
$\nabla^\dagger\nabla$ restricted to the imaginary quaternion sector
produces exactly the factor $3$ in the denominator.  No such calculation
exists.
\end{proof}

%──────────────────────────────────────────────────────────────────────────────
\section{Answers to the must-answer questions}
\label{sec:answers}
%──────────────────────────────────────────────────────────────────────────────

\begin{enumerate}[label=Q\arabic*.]

\item \textbf{Does $S[\Theta]$ naturally define a modular $\tau$ variable?}

  \textbf{Answer: Partially, with obstruction.}
  Complex time $\tau=t+i\psi$ is formally a complex variable residing in
  (a subset of) $\CC$.  If the imaginary-time direction is compactified
  with period $\beta$, the ratio $\tau/\beta$ can be treated as a
  complex-torus modulus in $\bH$.  However, $S[\Theta]$ is
  \emph{not} demonstrated to be invariant under $\SL(2,\ZZ)$ acting on
  this modulus (Obstruction O1).  The kinetic operator $\nabla^\dagger\nabla$
  has not been shown to transform as a modular covariant operator of
  definite weight.  Therefore the modular $\tau$ variable is
  \emph{formally available} but \emph{not physically justified}.
  \hfill\COND

\item \textbf{Does the $p$-winding sector transform under $\Gam_0(p)$?}

  \textbf{Answer: Arithmetically yes; physically unproved.}
  Definition~\ref{def:Hp} identifies $\Hp$ as the subspace of winding
  modes with $n\equiv 0\pmod{p}$.  By Lemma~\ref{lem:arith}, this
  subspace is closed under $\Gam_0(p)$ if $\SL(2,\ZZ)$ acts on winding
  modes.  But the physical action of $\SL(2,\ZZ)$ on winding modes
  has not been derived (Obstruction O2).  The arithmetic compatibility
  is exact; the physics derivation is missing.
  \hfill\COND

\item \textbf{Are there exactly $p+1$ inequivalent saddles?}

  \textbf{Answer: Arithmetically yes; physically unproved.}
  Lemma~\ref{lem:cosets} gives $|\Gam_0(p)\backslash\SL(2,\ZZ)|=p+1$
  exactly.  If each coset labels an inequivalent boundary condition for
  the path integral, then there are exactly $p+1$ candidate saddles.
  But their inequivalence in the UBT path integral (i.e., that no two
  saddles are gauge-related) has not been proved.
  \hfill\COND

\item \textbf{Are they equal-action?}

  \textbf{Answer: No — not derivable without resolving O1.}
  By Proposition~\ref{prop:equal}, equal action requires $\SL(2,\ZZ)$
  invariance of $S[\Theta]$, which is Obstruction O1.  Without O1,
  the $p+1$ saddles can have different classical actions and different
  one-loop determinants, and the degeneracy factor $p+1$ does not emerge.
  \hfill\NOGO

\item \textbf{Does their contribution multiply the $n\ln n$ coefficient?}

  \textbf{Answer: Conditionally yes.}
  Proposition~\ref{prop:nlogn} shows that, conditional on equal weighting
  (which requires O1--O3 to be resolved), the saddle sum $\Zp=(p+1)\cdot(\ldots)$
  contributes a factor $p+1$ to the one-loop free energy, which after
  taking $\ln$ produces a term $\sim(p+1)\ln p$ at $n=p$.  This matches
  the $B\cdot n\ln n$ structure in $V_{\mathrm{eff}}$ with $B=p+1$ before
  normalisation.  The normalisation then gives $B=(p+1)/3$.
  \hfill\COND

\item \textbf{Does the division by $3$ enter through modular area normalisation?}

  \textbf{Answer: Yes — this is an exact, unconditional statement.}
  Corollary~\ref{cor:area} and equation~\eqref{eq:ratio} establish that
  $B(p)=(p+1)/3=\vol(X_0(p))/\pi$, which is precisely the ratio of the
  level-$p$ modular curve area to $\pi$.  This follows from the exact
  identity $\vol(\SL(2,\ZZ)\backslash\bH)=\pi/3$ and is independent of
  all physics.  The physical question is only whether the path-integral
  normalisation over modular space is the correct one; this is reasonable
  but not derived from $S[\Theta]$.
  \hfill\PROVED\ (arithmetic); \COND\ (physics normalisation)

\end{enumerate}

%──────────────────────────────────────────────────────────────────────────────
\section{Summary of obstructions}
\label{sec:obstructions}
%──────────────────────────────────────────────────────────────────────────────

The construction fails at three successive levels.  Each obstruction is
sharp and precisely located.

\begin{center}
\renewcommand{\arraystretch}{1.4}
\begin{tabular}{@{}p{0.08\linewidth}p{0.36\linewidth}p{0.46\linewidth}@{}}
\toprule
Label & Obstruction & What would resolve it \\
\midrule
O1 & $S[\Theta]$ is not proved modular-invariant under $\SL(2,\ZZ)$ &
  Prove that $\nabla^\dagger\nabla$ on the complex-time torus is a
  modular-covariant operator of definite weight, and that the measure
  transforms covariantly. \\
\addlinespace
O2 & The $\SL(2,\ZZ)$-action on winding modes $\Theta_n$ has not been
  derived & Compute the action of a Möbius transformation on the Fourier
  expansion~\eqref{eq:winding} and show it maps $\Hp$ to itself
  (building on Lemma~\ref{lem:arith}). \\
\addlinespace
O3 & Equal-action of the $p+1$ saddles cannot be established without O1 &
  Resolve O1; then Proposition~\ref{prop:equal} gives equal action
  automatically. \\
\bottomrule
\end{tabular}
\end{center}

\begin{warnbox}
\textbf{Dependency.}  O3 depends on O1.  O2 is partially independent
and could in principle be resolved analytically without resolving O1.
Resolution of O1 and O2 together would collapse O3 and complete the
construction.
\end{warnbox}

%──────────────────────────────────────────────────────────────────────────────
\section{Synthesis table}
\label{sec:synthesis}
%──────────────────────────────────────────────────────────────────────────────

\begin{center}
\renewcommand{\arraystretch}{1.4}
\begin{tabular}{@{}p{0.44\linewidth}p{0.18\linewidth}p{0.28\linewidth}@{}}
\toprule
Claim & Status & Basis \\
\midrule
$|\Gam_0(p)\backslash\SL(2,\ZZ)|=p+1$ &
  \PROVED &
  Lemma~\ref{lem:index}--\ref{lem:cosets} \\
\addlinespace
$\vol(\SL(2,\ZZ)\backslash\bH)=\pi/3$ &
  \PROVED &
  Classical (Lemma~\ref{lem:fundvol}) \\
\addlinespace
$B(p)=\vol(X_0(p))/\pi=(p+1)/3$ (arithmetic identity) &
  \PROVED &
  Corollary~\ref{cor:area} \\
\addlinespace
$S[\Theta]$ defines a modular $\tau$ &
  \COND &
  O1 unresolved \\
\addlinespace
$\Hp$ transforms under $\Gam_0(p)$ &
  \COND &
  O2 unresolved \\
\addlinespace
$p+1$ inequivalent, equal-action saddles &
  \COND &
  O1, O2, O3 unresolved \\
\addlinespace
Saddle sum multiplies $n\ln n$ by $p+1$ &
  \COND &
  Conditional on all saddles equal-action \\
\addlinespace
Division by $3$ from modular area &
  \PROVED\ (arith.) \COND\ (physics) &
  Corollary~\ref{cor:area};  path-integral normalisation unproved \\
\addlinespace
$B(p)=(p+1)/3$ derived from $S[\Theta]$ &
  \NOGO (current) &
  O1--O3 block the derivation \\
\bottomrule
\end{tabular}
\end{center}

%──────────────────────────────────────────────────────────────────────────────
\section{Verdict}
\label{sec:verdict}
%──────────────────────────────────────────────────────────────────────────────

\begin{warnbox}
\textbf{Construction verdict: INCOMPLETE / NO-GO at current level.}

The arithmetic foundation is exact: the coset count is $p+1$, the
modular-area ratio gives $(p+1)/3$, and $\PP^1(\FF_p)$ provides the
correct index set for the saddles.

The physics construction fails at Obstruction O1: the action $S[\Theta]$
is not demonstrated to be invariant under $\SL(2,\ZZ)$ acting on the
complex-time modulus.  Without this invariance, the $p+1$ candidate
saddles cannot be shown to be equal-action, and the degeneracy factor
$p+1$ does not follow from the path integral.

Therefore $B(p)=(p+1)/3$ \textbf{remains a conditional modular ansatz}.
It is structurally coherent and arithmetically exact, but it is not a
consequence of $S[\Theta]$ until O1 and O2 are resolved.
\end{warnbox}

%──────────────────────────────────────────────────────────────────────────────
\section{Directions for resolution}
\label{sec:resolution}
%──────────────────────────────────────────────────────────────────────────────

\begin{enumerate}

\item \textbf{Prove modular invariance of $S[\Theta]$ on the torus.}
  The strategy would be to show that the biquaternionic Dirac--Laplace
  operator $\nabla^\dagger\nabla$ restricted to the complex-time torus
  $\CC/(\ZZ+\ZZ\tau)$ is a holomorphic family of operators parametrised by
  $\tau\in\bH$ and that its determinant transforms as a modular form of
  weight zero.  This is analogous to the modular properties of the
  string-theory one-loop amplitude and may admit a similar proof using
  $\eta$-function regularisation.

\item \textbf{Compute $\SL(2,\ZZ)$ action on winding modes.}
  Expand $\Theta(x,\tau)=\sum_n\Theta_n(x)\,e^{2\pi in\tau}$ in the new
  Fourier basis after $\tau\to\gamma\tau$.  Verify that the transformation
  matrices are block-diagonal on $\Hp$, so that $\Hp$ is $\Gam_0(p)$-stable.

\item \textbf{Check equal action numerically for small $p$.}
  For $p=2,3,5$, compute $S[\bar\Theta_a]$ explicitly for each of the
  $p+1$ coset representatives and check whether equality holds to leading
  order in a semi-classical expansion.  Failure would give a
  counter-example to the whole program.

\end{enumerate}

%──────────────────────────────────────────────────────────────────────────────
\section*{Internal references}
%──────────────────────────────────────────────────────────────────────────────

\begin{itemize}
  \item \texttt{research\_tracks/alpha\_spectral/hecke\_trace\_B\_derivation\_attempt.tex}
        — predecessor note; identifies Gaps H1--H3.
  \item \texttt{research\_tracks/alpha\_spectral/b\_coefficient\_gap\_resolution.tex}
        — synthesis of all routes to $B(p)$.
  \item \texttt{reports/hecke\_trace\_B\_verdict.md} — verdict on G137-B.
  \item \texttt{canonical/alpha/modular\_prime\_attractor\_theorem.tex}
        — formal statement of the prime-attractor theorem.
  \item \texttt{reports/hecke\_path\_integral\_no\_go\_or\_success.md}
        — companion report (this task's second deliverable).
\end{itemize}

\begin{thebibliography}{9}
\bibitem{Serre:1973}
J.-P.\ Serre, \textit{A Course in Arithmetic}, Springer, 1973.

\bibitem{Diamond:2005}
F.\ Diamond and J.\ Shurman, \textit{A First Course in Modular Forms},
Springer GTM 228, 2005.

\bibitem{Shimura:1971}
G.\ Shimura, \textit{Introduction to the Arithmetic Theory of Automorphic
Functions}, Princeton University Press, 1971.

\bibitem{Polchinski:1998}
J.\ Polchinski, \textit{String Theory, Vol.~1}, Cambridge University
Press, 1998.  (Modular invariance of the one-loop amplitude, Ch.~7.)
\end{thebibliography}

\end{document}
